\documentclass[11pt,a4paper]{article}
\usepackage[margin=1in]{geometry}
\usepackage{amsmath,booktabs,hyperref,float}

\title{Iteration 024: Mixup Augmentation + InstanceNorm}
\author{Autoresearch Loop}
\date{April 8, 2026}

\begin{document}
\maketitle

\begin{abstract}
Mixup augmentation ($\lambda \sim \text{Beta}(0.4, 0.4)$) with InstanceNorm
achieves $r=0.375$, not improving over the baseline ($r=0.378$). The
InstanceNorm component masks any potential benefit from mixup. Iter030
re-tests mixup without InstanceNorm to isolate its effect.
\end{abstract}

\section{Hypothesis}
Mixup creates virtual intermediate subjects by interpolating training windows:
$\tilde{x} = \lambda x_i + (1-\lambda) x_j$, $\tilde{y} = \lambda y_i + (1-\lambda) y_j$.
Since the mapping is approximately linear ($y \approx Wx$), mixup is theoretically
well-motivated. Combined with InstanceNorm for amplitude invariance.

\section{Method}
During training, each batch is mixed with a random permutation using
$\lambda \sim \text{Beta}(0.4, 0.4)$. Architecture: InstanceNorm + FIR (7 taps).
Combined loss + corr validation + 15\% channel dropout.

\section{Results}
\begin{table}[H]
\centering
\begin{tabular}{lrrr}
\toprule
Model & Mean $r$ & Std $r$ & SNR (dB) \\
\midrule
Combined loss (iter017) & 0.378 & 0.078 & 0.61 \\
Mixup + InstanceNorm & 0.375 & 0.076 & 0.61 \\
\bottomrule
\end{tabular}
\end{table}

\section{Analysis}
The result is confounded by InstanceNorm, which consistently degrades $r$
by $\sim$0.003. Mixup itself may be neutral or slightly positive, but we
cannot determine this from this experiment alone. Iter030 tests mixup
without InstanceNorm for a clean comparison.

\section{Next}
Re-test promising techniques (mixup, cosine loss, noise) without InstanceNorm.

\end{document}
